% ===========================================================================
%  Standalone front matter: title, abstract, conclusion.
%  This extension is written as its own paper; the closed-set study is cited
%  as prior work rather than assumed. Key for it: \cite{kbaier_closedset}.
% ===========================================================================

\title[mode=title]{The Threshold Does Not Transfer: Open-Set Test-Time Adaptation
Across Buildings}

\begin{abstract}
Deploying an activity recogniser in a new building means facing two shifts at once: the
sensors are in different places, and the occupant may do things the model was never
trained on. Open-set and universal domain adaptation address the second, and test-time
adaptation the first, but the methods for both are evaluated offline, with repeated passes
over a target set that is also the evaluation set. We study the combined shift under
constraints a deployment can meet: a source model trained once and frozen, a single online
pass over the unlabelled target stream in recording order, no source data, and no target
labels anywhere, including for the threshold that decides what counts as unknown.

Measured against a model trained on the target building at matched sample size, the
cross-building penalty on known activities is $1.6$ macro-F\textsubscript{1} points
entering one apartment and $36.0$ entering the other, a wider asymmetry than the
$0.6$ and $21.3$ of the closed-set case \cite{kbaier_closedset}. Four test-time methods
follow it: none improves on a frozen source model where there is no shift to recover.
One appears to, and the appearance is instructive. PETAL records the best open-set
H-score in one direction while classifying known activities roughly half as well, because
degrading a model uniformly lowers its confidence everywhere, pushes more windows below a
fixed threshold, and raises the unknown-detection term that H is bottlenecked by. The
standard open-set score is therefore maximised in part by damage, and pairing it with
known-class accuracy does not remove the problem; it only makes it visible.

The rejection rule fails for a different reason. Maximum softmax, entropy, maximum logit
and free energy are all at or below chance as novelty detectors here, with bootstrap
intervals excluding $0.5$ for three of the four: the model is as confident on activities it
never saw as on those it was trained on. A Mahalanobis distance to source class prototypes
is the only score above chance, and at a matched rejection rate it catches $1.58$ and
$1.17$ times as many unknowns as rejecting at random --- real, and far smaller than the
H-score comparison suggests, because two scores held at the same nominal percentile reject
very different fractions of a target stream.

That last observation is the main result. A threshold set to reject $1\%$ of source
validation rejects between $7\%$ and $63\%$ of the target stream, depending only on which
score carries it. Improving detection does not fix it: a source-only weighting rule that
gains $0.064$ AUROC in the direction with headroom \emph{loses} $0.11$ H at the threshold a
deployment would use, while the variant with the best H in both directions has
significantly worse AUROC in one. Detection quality and deployed open-set performance come
apart, so cross-domain threshold calibration is the binding constraint rather than one
difficulty among several. We recommend reporting the rejection rate, which is observable
without labels, beside every open-set score.
\end{abstract}

% ---------------------------------------------------------------------------
\section{Conclusion}
\label{sec:os-conclusion}

Adding a label-space shift to a geographic one makes cross-building activity recognition
harder without giving adaptation more to work with. Across a $2\times2$ grid of transfer
direction and label space, the methods track the measured penalty rather than the task:
they help only where there is a real shift, and in the direction where the penalty is
$1.6$ points none of them improves on doing nothing.

What survives is small and does not adapt. A frozen source model with a Mahalanobis
rejection rule, with prototypes and threshold fixed on source data before the target
stream is read, reaches H-scores of $0.5176$ and $0.4289$; every component we removed to
reach it --- the four adaptation methods, the confidence-based rejection scores, and the
learned embedding, which a distance on the raw thirteen features matches --- was worth
removing.

Three findings seem to us likely to outlive this dataset, and none of them is about our
system. The standard open-set score can be raised by degrading a model
uniformly whenever unknown detection is the weaker of its two terms, as it was in both
directions here, so it should never be reported without its known-class component. Four standard confidence-based
novelty scores are at or below chance here, which was visible only because a chance
baseline was computed, and is a reason to compute one before trusting any detector in a
new domain. And a threshold calibrated on source data did not transfer in either
direction here: the same nominal percentile rejects $7\%$ or $63\%$ of a target stream
depending on the score it sits on, and improving the detector underneath did not help
while that was true. We expect the mechanism to be general, since a threshold is a
statistic of the source distribution, but we have measured it on two buildings.

The limits are those of the study it extends. Two apartments of one dataset, one choice of
withheld activities, and a temporal axis that turns out to carry much of the novelty
signal in these particular homes. The specific numbers belong to that pair. What we would
expect to transfer is the procedure: measure the penalty per cell against a target-trained
model at matched sample size before running any method, compute a chance baseline for any
novelty score, report the rejection rate beside the H-score, and treat a method that wins
on H while losing on known-class accuracy as a null. Calibrating a rejection threshold
that holds across domains rather than across resamples of one is, on this evidence, the
problem worth working on next.
